\documentclass[header=false]{grad-statement}

\begin{document}

\section{Background}

As a teen, I spent five years in an accelerated math program at the University of Minnesota that covered the core of a typical undergraduate math degree.
Its rigor stood in sharp contrast to the classroom experience at my public high school where two-thirds of students failed to meet state reading standards and three-quarters received federal free or reduced-price lunch.
I fell in love with math and entered college thinking I'd be a mathematician.

But after taking several upper-level math classes, I became disillusioned with the idea of spending my life proving theorems understandable only by a handful of specialists.
Around the same time, a former high school track teammate was shot to death, and another close high school friend became a teen mother.
I began to question the social forces that had set us on such different trajectories, but the methods taught in introductory social science classes felt underpowered to grapple with such complicated issues.
I faced a conundrum:
Fields I found rigorous studied problems I found unsatisfyingly abstract.
Fields studying problems that moved me lacked sufficient rigor.

Gradually, I realized that the choice between rigor and relevance was a false dichotomy.
Since college, I've had the privilege of working at Opportunity Insights and tackling issues I find personally meaningful---education, social mobility, health inequality---with the exactness that first drew me to mathematics.
My individual work and the research environment have been intellectually challenging (how do we efficiently bootstrap millions of standard errors?) and personally fulfilling (how do we develop policies that enable children to rise from poverty?).
The IDSS PhD program is an ideal setting to continue doing work that checks both these boxes.

\section{Interests}

% (500 words max) What kinds of research connected to data, systems, and society interest you? Are there particular issues and problems you wish to address? Include, as far as you can, your particular interests, whether experimental, theoretical, or application oriented. Explore the intersection of your interests, your preparation, and MIT/ IDSS programs and research. Feel free to mention specific research and researchers. Note: It is OK if your interests are not yet settled. If this is the case, use this space to explore one or two areas of overlap between your interests and IDSS-related work.

My experience growing up in an urban public school system and my recent work at Opportunity Insights (OI) have made me particularly interested in the unequal distribution of opportunity across America.
Why are some children---due simply to where they are born---more likely to succeed in adulthood than others?
At OI, I've worked on projects that descriptively document the extent of this inequality and projects that look for ways to level the playing field and give everyone the chance to rise out of poverty and achieve the American Dream.
%
In graduate school, I'd like to continue studying this topic, but from a slightly different angle:
What are the downstream social consequences when opportunity disappears?

For example, in recent independent work, I've developed a measure of ideological polarization in the American public based on the eigenvalue spectrum of the covariance matrix of people's opinions.
Intuitively, polarization is maximized when people disagree strongly on individual topics and when that disagreement is highly correlated across issues.
I show that this index has been increasing over the last several decades, and this increase is not explained by the divergence of Republican and Democratic opinions. In fact, much of this increase has occurred within political party.

I'd like to continue this work in graduate school by connecting it to the decline of opportunity:
Has polarization increased in places where opportunity has declined?
To what extent is rising polarization due to diverging opinions of those in places where opportunity remains alive vs those in places where it has disappeared?
Outside this particular idea, I hope to pursue a broader research agenda about spatial inequality in the United States from a variety of angles corresponding to my diverse background in mathematics, statistics, computer science, and economics.
I've appreciated how my current work at Opportunity Insights is willing to use whatever technique is necessary to bring novel quantitative clarity to publicly relevant issues.
I hope to continue this in my own research.

Outside this applied work, I'm also interested in pursuing more technical research in algorithms, statistics, and machine learning.
I expect that my interests will continue to evolve over the course of graduate school, and the SES PhD provides a unique combination of breadth to allow this exploration and deep faculty expertise to enable meaningful research in whatever area I eventually settle.

\section{Purpose}

% (300 words max) Discuss your motivation for pursuing a PhD in Social and Engineering Systems. Suggested prompts: What are your long term goals after graduation? What types of contributions do you see yourself making in your field, to society, etc.? If you are still deciding, what kinds of work are potentially interesting, valuable, or meaningful to you?

After my PhD, I hope to become an academic conducting socially-relevant research at the intersection of economics, statistics, and computer science.
In my work at Opportunity Insights, I have seen how academic research---when put in the correct hands---can lead to meaningful change in the real world.
The lab's work is frequently covered by major news outlets and is widely cited in policy discussions including the President's State of the Union address.
During and after my PhD, I hope to produce research of similar societal value.

Outside my research, I also plan to have teaching and mentorship remain a core part of my career.
At Carleton College, I served as a TA in eight of my eleven terms on campus.
I loved sharing my excitement with younger students and encouraging those who had previously been dissuaded from math and computer science by their reputation as overly technical and disconnected from reality.
In my current job as a research assistant for the Opportunity Insights group at Harvard, I've continued this interest in teaching and mentorship by helping spearhead the development of training materials and leading several Python and Git crash courses to help our new hires learn the coding skills necessary for empirical economics research.
I plan to continue honing my teaching skills during graduate school by serving as a TA as often as I'm able.

Throughout my life, I've valued leaders and teachers whose infectious enthusiasm lit a fire under me and made me achieve things that I thought I wasn't capable of.
I hope to provide this same mentorship to the next generation of students during graduate school and throughout my career as a professor.

\section{Additional Info}

I have a strong commitment to producing public goods in the form of software.

In my work at OI, I've often been frustrated by how programming languages that provide powerful statistical modeling---like Stata and R---are painfully slow on big datasets and awkward to use when managing large complex codebases.
Meanwhile, tools that work well at scale---like Python and its excellent Polars data frame library---lack good statistical modeling software.
Switching between languages creates an annoying friction point and increases the chances of bugs.

Over the past year, I've worked to fill this gap by writing a statistical library for Polars on par with the capabilities of R and Stata.
So far, I've implemented features including OLS and generalized linear models that allow for sample weights, regularization, several spline bases, and common sandwich covariance estimators.
I'm currently working to add random effects, REML hyper-parameter selection, and the ability to absorb high-dimensional fixed effects.
I plan to continue developing this library throughout graduate school and hope that it soon removes many of the headaches that come with a complicated empirical research project.

The Interdisciplinary Doctoral Program in Statistics offered by IDSS provides a unique opportunity to combine this interest in statistical software with my broader interest in impactful social research.

\end{document}
